\section{Conclusion}
\label{sec:conclusion}

In this project we studied the stochastic shortest path problem through the lens of portfolio optimization. We framed the problem as selecting, for each source--target pair, a small set of candidate paths in a preprocessing phase, such that at query time we can return the cheapest among them under the realised edge costs. We implemented three portfolio algorithms---a greedy selection over the full path set, a frequency-based selection of the paths that are shortest the most often, and a penalty-based scheme that iteratively penalises edges of expected-cost shortest paths---and evaluated them on a real-world traffic dataset from Shanghai.

Our main finding is that on this instance, portfolios of moderate size achieve near-optimal expected costs. The greedy and frequency-based algorithms both reach expected cost increases within hundredths of a percent compared to the optimal with $k \approx 8$ paths, and the frequency-based approach does so without enumerating the exponentially many candidate paths. This suggests that for traffic networks with similar characteristics, the portfolio model can provide a meaningful reduction in query-time computation without a substantial loss in solution quality. The penalty-based scheme, while cheaper to precompute, converges more slowly, indicating that information beyond the mean edge cost matters when building diverse portfolios.

These results come with several caveats. The Shanghai dataset is small ($n = 67$, $m = 156$) and covers a limited geographic area where trips are short. The graph is a subgraph of a larger network and may contain artifacts from the way it was extracted. It is not clear how well the findings transfer to larger graphs, longer trips, or different types of traffic data. Our evaluation is also limited to one instance and one distribution; the difficulty of the stochastic shortest path problem depends heavily on the graph structure and the nature of the cost uncertainty, as we discussed in \cref{sec:stats}.

\subsection{Future work}

A natural next step is to evaluate the portfolio approach on additional datasets. The Shanghai dataset used here is small; datasets for Shenzhen ($156$ vertices) and Los Angeles ($207$ vertices) are available from the other sources and would provide a useful check on whether the trends we observed persist in larger graphs with different connectivity patterns. But it is difficult to obtain continuous time series and for example the data from Shenzhen has time intervals with no recordings. There is a point to be made towards defining the cost distribution differently, instead of a uniformly chosen timepoint throughout the month, one might want to sample uniformly from all observed car speeds, giving much more weight to rush hours and eliminating the problem of time intervals without traffic. Datasets with longer time spans or coarser time aggregation might also change the nature of the uncertainty and thus the value of portfolios.

On the theoretical side, it would be interesting to study if one can identify conditions on the distribution $\mathcal D$ under which a constant-factor approximation guarantee can be proved for one of the proposed algorithms. For example, if the edge cost distributions are independent and have bounded variance, one might be able to bound how far the expected-cost shortest path can be from the true shortest path in any realisation. Such results would complement the empirical evaluation we presented here.

The memory cost of storing per-pair portfolios grows as $\mathcal O(k n^3)$, which becomes prohibitive for graphs with thousands of vertices. The unfinished work in \cref{sec:graphwide} points toward data structures that store a single graph-wide set of paths and derive per-pair portfolios from it on the fly. There are also connections to learning-based approaches: one could train a model to predict, from some features of the query and the current traffic conditions, which of a small set of global routes is likely to be cheapest. Such methods could trade some theoretical guarantees for better scalability.
